\documentclass[11pt,a4paper]{article}
\usepackage[margin=25mm]{geometry}
\usepackage[T1]{fontenc}
\usepackage{lmodern,amsmath,amssymb,amsthm,booktabs,longtable,array,microtype,enumitem,xcolor}
\usepackage[colorlinks=true,linkcolor=blue!50!black,urlcolor=blue!50!black,citecolor=blue!50!black]{hyperref}
\usepackage{fancyhdr}
\pagestyle{fancy}\fancyhf{}\fancyhead[L]{\small Astra P1: one feedback round}\fancyhead[R]{\small 31 October 2026, Brisbane}\fancyfoot[C]{\thepage}
\setlength{\headheight}{14pt}\setlength{\parindent}{0pt}\setlength{\parskip}{6pt}
\setlist{nosep,leftmargin=1.5em}
\newcommand{\E}{\mathbb E}\newcommand{\Var}{\operatorname{Var}}
\newcommand{\tr}{\operatorname{tr}}\newcommand{\rank}{\operatorname{rank}}
\newcommand{\norm}[1]{\left\lVert#1\right\rVert}
\newcommand{\Rad}{R^*_{\mathrm{ad}}}
\newtheorem{theorem}{Theorem}\newtheorem{lemma}[theorem]{Lemma}
\newtheorem{proposition}[theorem]{Proposition}\newtheorem{corollary}[theorem]{Corollary}
\title{\textbf{One feedback round for near-optimal quadratic measurement}\\[5pt]
\large Signed spectral compression, an all-noisy separation, and acquired implementations}
\author{Astra P1: central construction with independent internal challenges}
\date{31 October 2026, Australia/Brisbane}
\hypersetup{pdftitle={One feedback round for near-optimal quadratic measurement},pdfauthor={Astra P1},pdfsubject={Quadratic functional prediction, adaptive linear measurements, signed spectral bounds, noise and acquisition}}
\begin{document}\maketitle
\begin{abstract}
For a known quadratic target $f(X)=X^TAX+b^TX+c$ under $X\sim N(0,I)$, an explicit two-round experiment uses at most $k+1$ exact scalar linear observations and has residual squared risk no greater than every unrestricted adaptive $k$-observation policy. The first batch pairs positive and negative eigendirections so their unobserved quadratic compression vanishes; one final observation reads the remaining outcome-dependent linear term. Compression interlacing supplies a matching comparison lower bound, and the expected-budget version follows by convex interpolation. Exact worst-case query complexity is classified by inertia. A separate all-noisy theorem for a normalized Gaussian inner product proves a nonvanishing accuracy advantage against the optimal fixed experiment, including repeated and nonorthogonal directions: at dimension 256 and noise variance 0.01, 129 adaptive reads have risk 0.019704 versus optimal fixed risk 0.506023; 256 fixed reads match the adaptive risk. We construct a bounded-error quadratic acquisition route, prove model-error transfer, execute a complete synthetic acquisition/planning/deployment fixture, and retain the costs of dense directions and physical measurement access. Coordinate-only and grouped-acquisition companions are preserved separately. Historical priority and practical device advantage are unresolved.
\end{abstract}

\section{The capability and the changed question}

The preceding P1 theorem replaced arbitrary adaptive coordinate planning by a larger fixed set, under low interaction order and controlled coordinate fluctuation moments. It preserved Bayesian residual risk, with a conservative measurement multiplier. The rare-gate family showed why feedback can find valuable observations that are not valuable unconditionally. Those results, acquisition constructions, corrections and scoped reviews remain preserved.

The present question is whether a much smaller amount of feedback can retain the useful part of adaptive experimentation. The mechanism is to choose a first batch that simplifies the \emph{unobserved target}, then acquire its remaining linear functional in one feedback step. The first batch need not maximize immediate explained variance. It may deliberately retain uncertainty whose influence will become observable after the batch.

One relevant scientific objective is goal-oriented sensing: collect data to reduce uncertainty in a nonlinear quantity of interest. Quadratic approximations and their exact Gaussian variance objective are already developed in the primary work of Neuberger, Alexanderian and van Bloemen Waanders \cite{goal}. Their fixed sensor-placement framework and matrix-free computations are strong established alternatives. The present question concerns the additional capability of choosing a linear measurement after seeing a first batch. Arbitrary programmable directions are a stronger access model than a supplied menu of sensor locations; a transfer to such applications must earn that access and the quadratic approximation error.

For the central result, a measurement is a scalar linear functional of the unknown input. This is a substantive change from coordinate-only access. All compared methods receive the same permitted measurements, known input law and target information. If forming a dense linear measurement requires individually acquiring every coordinate, its cost must include those acquisitions. The theorem does not turn a coordinate sensor into a cheap aggregate sensor.

The report establishes mathematical and executable capabilities under the stated interfaces. It does not establish an unknown Gaussian prior, certify a quadratic model class from finite agreement, or claim a speed advantage from scalar-query counts alone. The known ingredients are Gaussian conditioning, interlacing, quadratic-form inertia, polynomial interpolation and elementary statistics. The particular combined information-complexity statements were derived here and independently reconstructed; historical priority remains open.

\section{Exact linear measurements and the signed spectral bound}

Let $X\sim N(0,I_n)$ and $A=A^T$. A policy chooses a scalar linear observation from its preceding outcomes and independent private randomness, and may stop. Remove redundant observations and orthogonalize each new direction against earlier ones. An exact dependent observation contains no new information. Let $K$ count nonredundant scalar reads and let $\mathcal H$ contain the full transcript and independent seed.

The Bayes risk is $R_\pi=\E\Var(f(X)\mid\mathcal H_\pi)$. We use either the hard constraint $K\le k$ or the expected constraint $\E K\le k$, and state which one is intended. All policies are measurable and return almost surely. The conditional-mean estimator is the best squared-loss estimator for their transcript.

List the positive eigenvalues as $\lambda_1\ge\cdots\ge\lambda_p>0$ and magnitudes of negative eigenvalues as $\mu_1\ge\cdots\ge\mu_q>0$. Define
\[
 \mathcal B(k)=2\left(\sum_{j>k}\lambda_j^2+\sum_{j>k}\mu_j^2\right)
 \quad\text{for integer }k\ge0. \tag{1}
\]
For real $k\ge0$, linearly interpolate between adjacent integer values. This is decreasing and convex, because its successive decrements are $2(\lambda_{k+1}^2+\mu_{k+1}^2)$, with absent eigenvalues taken as zero.

\begin{lemma}[Gaussian transcript and residual risk]\label{lem:conditional}
For almost every adaptive transcript, the conditional law is $X=m+Z$, where $m$ is the observed orthogonal projection, $Z\sim N(0,Q)$, and $Q$ is the projection onto the unobserved complement. At that transcript,
\[
 \Var(f\mid\mathcal H)
 =2\norm{QAQ}_F^2+\norm{Q(b+2Am)}^2. \tag{2}
\]
\end{lemma}
\begin{proof}
Condition on the independent seed. If the next unit direction $u$ is chosen in the current unobserved space, its innovation $u^TZ$ is standard Gaussian and the remaining orthogonal Gaussian component is independent of it. Induction proves the conditional law even when the chosen directions and stopping rule depend on earlier outcomes. On a fixed transcript, $m,Q$ are fixed. Expand the target into its known constant, the linear term $[Q(b+2Am)]^TZ$, and $Z^TQAQZ$. The centered linear and quadratic Gaussian terms are orthogonal, and the quadratic variance is $2\norm{QAQ}_F^2$.
\end{proof}

\begin{lemma}[Signed tail lower bound]\label{lem:lower}
Every hard-$k$ adaptive policy has $R_\pi\ge\mathcal B(k)$. Every expected-budget-$k$ policy has the same bound with the interpolated definition (1).
\end{lemma}
\begin{proof}
Compression interlacing implies that a codimension-$t$ compression $QAQ$ retains positive eigenvalues at least $\lambda_{t+1},\lambda_{t+2},\ldots$, and negative eigenvalue magnitudes at least $\mu_{t+1},\mu_{t+2},\ldots$. Hence
\[
 2\norm{QAQ}_F^2\ge\mathcal B(t).
\]
Drop the nonnegative linear term in (2). The hard-budget assertion follows from monotonicity. For random $K$, take expectations and use Jensen:
$R_\pi\ge\E\mathcal B(K)\ge\mathcal B(\E K)\ge\mathcal B(k)$.
\end{proof}

The spectral minimum is classical matrix machinery. Fan--Pall's compression theorem \cite{fanpall} characterizes all codimension-$k$ spectra by interlacing; choosing each admissible compressed eigenvalue nearest zero gives exactly the minimum $2\norm{QAQ}_F^2=\mathcal B(k)$. Its real-symmetric version permits orthogonal changes of basis. The pairing below is an explicit construction of this minimum. The bound does not assert that the conditional linear term is absent. Removing that term constructively and comparing against adaptive information is the next step.

\section{The two-round compiler}

For a positive/negative eigenpair $Ae_+=\lambda e_+$ and $Ae_-=-\mu e_-$, define orthonormal vectors
\[
 u=\frac{\sqrt\lambda e_+-\sqrt\mu e_-}{\sqrt{\lambda+\mu}},
 \qquad
 v=\frac{\sqrt\mu e_++\sqrt\lambda e_-}{\sqrt{\lambda+\mu}}. \tag{3}
\]
In the coordinates $U=u^TX$, $V=v^TX$, that target block becomes
\[
 (\lambda-\mu)U^2+2\sqrt{\lambda\mu}\,UV. \tag{4}
\]
Observing $U$ leaves no quadratic term in $V$. Its unresolved effect is linear, with a coefficient determined by the observed $U$.

Pair the leading positive and negative modes up to $\min(k,p,q)$. If one sign runs out, use ordinary eigenvector reads for the remaining leading modes of the other sign, up to $k$ first-stage directions in total. Write their orthonormal matrix as $U_0$, put $P=U_0U_0^T$, and $Q=I-P$. The compression
\[
 D=QAQ
\]
contains the untouched signed spectral tails and zero paired blocks. Therefore
\[
 2\norm D_F^2=\mathcal B(k). \tag{5}
\]

After the first batch, compute $m=PX$ and
\[
 t=Q(b+2Am).
\]
If $t\ne0$, use the single final measurement $z=w^TX$ with $w=t/\norm t$. If $t=0$, omit this measurement. The final conditional-mean predictor is
\[
 \widehat f=m^TAm+b^Tm+c+\norm t\,z+\tr D
             +(w^TDw)(z^2-1). \tag{6}
\]
When the last measurement is omitted, omit its two terms as well.

\begin{theorem}[One extra observation replaces unrestricted sequential planning]\label{thm:main}
The above explicit two-round policy uses at most $\min(n,k+1)$ exact scalar reads and satisfies
\[
 \boxed{R_{\mathrm{two\ rounds}}(k+1)
          \le\mathcal B(k)\le\Rad(k).} \tag{7}
\]
Here the right comparator may be any hard-$k$ policy, or any expected-budget-$k$ policy for integer $k$. No comparator policy is supplied or computed. For real expected budget $k$, randomize the first-stage size between $\lfloor k\rfloor$ and $\lceil k\rceil$ with the interpolation probabilities. The resulting expected total budget is at most $k+1$ and its risk is at most $\mathcal B(k)\le\Rad(k)$.
\end{theorem}
\begin{proof}
Given the first batch, write $X=m+Z$, $Z\sim N(0,Q)$. The target is its known constant plus $t^TZ+Z^TDZ$. The final read makes the linear term known. Conditional Gaussian projection gives
\[
 \E[Z^TDZ\mid z]=\tr D+(w^TDw)(z^2-1),
\]
which proves (6). The conditional-first-stage risk, averaged over the last read and the unobserved Gaussian components, is exactly
\[
 2\norm D_F^2-2(w^TDw)^2\le 2\norm D_F^2. \tag{8}
\]
To see that no quadratic cross term is missing, put $R=Q-ww^T$. The risk at a completed last-read transcript is
$2\norm{RDR}_F^2+4z^2\norm{RDw}^2$.
Averaging $z^2=1$ and decomposing the Frobenius norm gives (8).
Equations (5), (8) and Lemma~\ref{lem:lower} prove (7). If $b=0$, the final direction is supported in the paired zero-compression space, so $Dw=0$ and the achieved risk is exactly $\mathcal B(k)$. The randomized expected-budget construction averages adjacent values of (5).
\end{proof}

The use of a supremum or infimum does not presume an optimizing policy exists: the lower bound holds for every comparator separately. The construction is spectral algebra plus one outcome-dependent vector, with no Bellman search over histories.

\subsection{Strong fixed designs and semidefinite targets}

For homogeneous $f=X^TAX+c$, the best fixed $r$-dimensional linear experiment has
\[
 R_{\mathrm{fixed}}^*(r)=2\sum_{j>r}s_j(A)^2, \tag{9}
\]
where $s_j(A)$ are absolute eigenvalues in decreasing order. Indeed, a fixed projection $P$ explains $2\norm{PAP}_F^2$; this is at most $2\norm{PA}_F^2$, whose maximum is the leading squared-singular-value sum. Its eigenvectors attain equality. This is the comparator used, rather than an uninformed random sketch.

For $A$ positive or negative semidefinite and $b=0$, (1) and (9) coincide at the same $k$. Adaptive exact linear measurements have no advantage in Bayes risk in this case. The benefit of pairing arises from opposing signs, not simply from the presence of a quadratic term.

For indefinite $A$, fixed leading modes of both signs use at most
$k+\min(k,\min(p,q))$ reads and attain $\mathcal B(k)$; include the direction $b$ at one additional read if needed. This records an ordinary fixed alternative using the same acquired spectrum.

\subsection{One extra read cannot be silently removed}

The explicit policy is not generally optimal at its own budget. For $f=2U^TV$ with independent $m$-dimensional standard Gaussian blocks and $m\ge2$, both the simple paired two-read policy and optimal fixed two-read experiment explain variance $4$. A better adaptive policy reads $U_1=y$, then reads $V_1$ when $4y^2\ge2$, and otherwise reads an unused eigenvector. Its value is
\[
 2+\E(4G^2-2)_+=4.716565402\ldots,\qquad G\sim N(0,1).
\]
The independent audit proves this is the full two-read optimum by reducing a first direction to its Rayleigh quotient and optimizing the second direction. This threshold effect is retained as a useful boundary: (7) is an additive-budget comparison, not a same-budget exact optimizer.

\section{Exact query complexity and the role of inertia}

The exact reconstruction construction is pointwise and does not require a probability law. Let $A$ have inertia $(p,q,z)$, meaning $p$ positive, $q$ negative and $z$ zero eigenvalues. A maximal totally $A$-isotropic subspace $L$ has dimension $z+\min(p,q)$. The paired vectors $v$ in (3), together with the nullspace, construct one. Reading $L^\perp$ costs $\max(p,q)$ observations. The target restricted to the remaining affine fiber is linear; one final direction reads it exactly.

\begin{theorem}[Worst-case exact linear-query classification]
For an arbitrary real quadratic target, the optimal hard number of exact adaptive scalar linear queries is
\[
 \begin{cases}
 0,& A=b=0,\\
 \rank A,& A\text{ is semidefinite and }b\in\operatorname{range}(A),\\
 \rank A+1,& A\text{ is semidefinite and }b\notin\operatorname{range}(A),\\
 \max(p,q)+1,& A\text{ is indefinite}.
 \end{cases} \tag{10}
\]
The same hard-count lower bound applies to zero-error almost-sure evaluation under a full-support Gaussian law. This statement does not classify the optimal expected exact count.
\end{theorem}
\begin{proof}
Every terminal transcript freezes its queried directions. Its fiber is an affine subspace. Constancy of $f$ on that fiber forces its direction $L$ to be totally $A$-isotropic and perpendicular to $2Ax+b$. Thus at least $\max(p,q)$ reads are necessary.

At exactly that count, $L$ is maximal isotropic and contains $\ker A$. If $b$ has a nullspace component, perpendicularity is impossible. Otherwise complete the square using $y=x+\tfrac12 A^\dagger b$. Suppose $p\ge q>0$. The form on $L^{\perp_A}$ is positive semidefinite, with radical $L$ and positive quotient dimension $p-q$. Constancy therefore cannot hold at any input satisfying $y^TAy<0$, a nonempty open set. For $p<q$, use the positive cone. This proves the extra read is necessary for indefinite $A$. The semidefinite cases follow because its maximal isotropic subspace is its nullspace; observe the range and, when required, the remaining linear direction. The constructed upper bounds attain all cases.

For the Gaussian almost-sure version, conditional fibers carry absolutely continuous Gaussian densities. A polynomial that is constant almost everywhere on such a fiber is constant on it. The obstructing open cone has positive probability. An independent random seed does not change this argument.
\end{proof}

In particular $f=2U^TV$ is computed exactly using $m+1$ adaptive reads, while every fixed exact experiment needs $2m$. The ratio tends to two. Exact-real coefficients and observations are part of (10); finite precision requires the risk certificates below.

\section{A complete noisy comparison at nonvanishing accuracy}

Let
\[
 F=\frac{U^TV}{\sqrt m},\qquad U,V\stackrel{\mathrm{ind}}\sim N(0,I_m),
 \qquad \Var(F)=1.
\]
Every measurement has a unit direction and independent Gaussian noise of variance $s>0$. All methods are allowed arbitrary unit directions in the same $2m$-dimensional space. Put $\alpha=(1+s)^{-1}$.

Read $Y=U+\eta$ in the first batch. Next read
\[
 Z=\frac{Y^TV}{\norm Y}+\xi,
 \qquad \widehat F=\frac{\alpha^2\norm Y\,Z}{\sqrt m}. \tag{11}
\]
The zero-norm event has probability zero; any defined fallback is allowed there.

\begin{theorem}[All-noisy two-round advantage against optimal fixed sensing]\label{thm:noisy}
The protocol (11) uses $m+1$ noisy reads and has Bayes risk
\[
 R_{\mathrm{two}}=1-\frac1{(1+s)^2}. \tag{12}
\]
For $0<s\le1$ and $r\le2m$, the optimal fixed experiment, including repeated and nonorthogonal unit directions and arbitrary estimators, has risk
\[
 \boxed{R^*_{\mathrm{fixed}}(r)=1-\frac{r}{2m(1+s)^2}.} \tag{13}
\]
Thus $2m$ fixed reads match the risk attained with $m+1$ adaptive reads.
\end{theorem}
\begin{proof}
Given $Y$, the posterior mean of $U$ is $\alpha Y$. The independently observed projection of $V$ has posterior mean $\alpha ZY/\norm Y$. Their conditional means multiply, giving (11). Its variance is $\alpha^2$, since $\E\norm Y^2=m(1+s)$ and $\E[Z^2\mid Y]=1+s$. Orthogonality gives (12).

For a fixed measurement matrix $M$ with $r$ unit rows, put $H=M^TM$ and $P=H(H+sI)^{-1}$. The posterior-mean covariance of the input is $P$. Represent the target as $X^TAX$ with
\[
 A=\frac1{2\sqrt m}\begin{pmatrix}0&I\\I&0\end{pmatrix},
 \qquad A^2=I/(4m).
\]
The explained variance is $2\tr(APAP)\le\tr(P^2)/(2m)$ by Frobenius Cauchy--Schwarz. Let $t_j$ be the eigenvalues of $H$, padded to $r$ values. Their sum is $r$. For $g_s(t)=t^2/(s+t)^2$, direct factorization gives
\[
 g_s(1)+g_s'(1)(t-1)-g_s(t)
 =\frac{s(t-1)^2(2t+s-s^2)}{(1+s)^3(s+t)^2}\ge0 \tag{14}
\]
when $t\ge0$ and $s\le1$. Summing implies $\tr(P^2)\le r\alpha^2$. Orthogonal eigenvector measurements attain equality. Randomized fixed designs obey the same bound after averaging.
\end{proof}

At $m=128,s=0.01$, the exact comparison is
\begin{center}
\begin{tabular}{lrr}\toprule
Experiment & Reads & Residual risk\\\midrule
Two rounds &129&0.0197039506\\
Optimal fixed, same budget &129&0.5060226939\\
Fixed matching the two-round risk &256&0.0197039506\\\bottomrule
\end{tabular}
\end{center}
This is a comparison at unit prior variance and nonzero noise on \emph{every} read. It is not a ratio obtained by dividing two vanishing explained variances.

The range $s\le1$ in the fixed-design optimum is substantive. At $s=2,m=1$, two repeated eigenvector reads explain $1/8$, while two orthogonal reads explain $1/9$. Equation (12) remains true at larger noise, but (13) cannot be extended unchanged.

\subsection{General quadratic noisy reconstruction of a compressed target}

The mechanism remains useful beyond equal paired eigenvalues. Suppose the first fixed batch has orthonormal matrix $U_0$ and makes $QAQ=0$. Write the observed block as $H=U_0^TAU_0$ and the cross block as $C=QAU_0$. With noisy first values $Y$, let $\mu=\alpha Y$, $v=1-\alpha$, and
\[
 C_0=c+\mu^TH\mu+v\tr H+b_U^T\mu,
 \qquad t=b_Q+2C\mu.
\]
One last noisy unit-direction read along $t$ has Bayes predictor $C_0+\alpha\norm t\,Z$. Its risk is exactly
\[
 2(1-\alpha^2)\norm A_F^2+(1-\alpha)\norm b^2. \tag{15}
\]
The first and final explained-variance contributions are, respectively,
$2\alpha^2\norm H_F^2+\alpha\norm{b_U}^2$ and
$4\alpha^2\norm C_F^2+\alpha\norm{b_Q}^2$.
Since $QAQ=0$, they sum to $2\alpha^2\norm A_F^2+\alpha\norm b^2$, proving (15).

This is the risk of observing a full orthonormal basis with the same noise, achieved in at most $\max(p,q)+1$ reads (or fewer when the final direction vanishes). A full noisy orthonormal experiment need not be the optimal fixed experiment for an arbitrary unequal spectrum. We do not claim a general noisy counterpart of Theorem~\ref{thm:main} against every noisy adaptive policy.

For a partial homogeneous paired design, the achieved explained variance is $\alpha^2$ times the selected positive-plus-negative spectral variance. The last read captures all selected bilinear mean terms together. Unselected spectral tails retain their prior uncertainty.

\subsection{Moment robustness, calibration and finite direction precision}

The \emph{achievability} calculation in (11) needs only independent centered isotropic blocks $U,V$ and independent zero-mean readout errors. Gaussianity is required for its Bayes interpretation and the optimal-fixed lower bound. If first- and second-stage noise variances are $s_1,s_2$, put $W=\norm Y\,Z/\sqrt m$. Then
\[
 \E(FW)=1,\qquad \E W^2=(1+s_1)(1+s_2).
\]
Using $\widehat F=\beta W$ with a certified upper bound $D\ge(1+s_1)(1+s_2)$ and $\beta=1/D$ gives risk at most $1-1/D$. Calibration is a required input; this bound does not certify noise independence or isotropy from a small pilot.

For a finite-precision first vector $\widetilde Y=Y+e(Y)$ with $\E\norm e^2/m\le\tau_Y^2$, coupled use of the same second noise gives an $L^2$ change in $W$ at most $\sqrt{1+s_2}\,\tau_Y$. The error must be determined from the first transcript and independent private randomness, before seeing $V$ or the last noise. An additional last-read error bounded by $\delta_Z$ contributes at most $\delta_Z(\sqrt{1+s_1}+\tau_Y)$. Multiply their sum by $\beta$ and add it to the square root of ideal risk. The direction $\widetilde Y/\norm{\widetilde Y}$ need not be well conditioned near zero, but the weighted target functional remains controlled by this bound. Programming a length-$m$ direction still costs $m$ coefficients and their precision bits. This is not a free one-bit feedback message.

\section{Acquiring the model and paying numerical work}

\subsection{An explicit scalar source construction}

For a genuinely promised quadratic function available at chosen real inputs, query $f(0)$, $f(e_i)$, $f(-e_i)$, and $f(e_i+e_j)$ for $i<j$. They give
\begin{align*}
 c&=f(0),\\
 b_i&=[f(e_i)-f(-e_i)]/2,\\
 A_{ii}&=[f(e_i)+f(-e_i)-2f(0)]/2,\\
 A_{ij}&=[f(e_i+e_j)-f(e_i)-f(e_j)+f(0)]/2.
\end{align*}
There are exactly $N=(n+1)(n+2)/2$ scalar evaluations. Reading an already supplied dense coefficient matrix costs $\Theta(n^2)$ instead. The scalar evaluations each pay the source's own execution cost, and a dense interface submits $nN$ coordinates. A sparse-input interface is a separate access contract.

If each scalar value has absolute error at most $\nu$, then $|\Delta c|\le\nu$, $|\Delta b_i|\le\nu$, and $|\Delta A_{ij}|\le2\nu$. Under the known standard Gaussian law,
\[
 \norm{f-\widehat f}_{L^2}^2
 =2\norm{\Delta A}_F^2+\norm{\Delta b}^2+
                 (\tr\Delta A+\Delta c)^2
 \le(12n^2+5n+1)\nu^2. \tag{16}
\]
The trace/mean error is included. If readout errors are Gaussian and averaged to establish simultaneous absolute bounds, those repeats and their confidence failure must be paid separately. A supplied bound $\nu$ is not a learned calibration result.

Finite stencil agreement does not establish quadratic adequacy. For example,
$h(x)=Lx_1(x_1^2-1)$ vanishes at all these probes but has Gaussian squared norm $10L^2$. The model promise or an independently justified approximation bound remains necessary.

\begin{proposition}[Model-error transfer]\label{prop:robust}
Under the same known input law and measurement interface, suppose
$\norm{f-\widehat f}_{L^2}\le\tau$. Compile the two-round policy for $\widehat f$ and deploy its own conditional predictor. Then
\[
 \boxed{\sqrt{\operatorname{MSE}_{f}}
 \le\sqrt{R^*_{\mathrm{ad},f}(k)}+2\tau.} \tag{17}
\]
\end{proposition}
\begin{proof}
For any transcript $\mathcal H$, conditional expectation and its complementary residual are $L^2$ contractions. Hence the square roots of the Bayes risks of $f$ and $\widehat f$ at the same policy differ by at most $\tau$, and so do their optimal square-root risks over a shared policy class. The actual error from the deployed $\widehat f$ predictor is at most $\sqrt{R_{\widehat f}(\mathcal H)}+\tau$ by the triangle inequality. Apply Theorem~\ref{thm:main} to $\widehat f$ and combine.
\end{proof}

For the noisy full-compression formula (15), the actual target's root MSE is at most the square root of its acquired-model formula plus $\tau$. This uses the same first triangle inequality and does not assert noisy optimal-adaptive equivalence.

\subsection{Spectral preparation and an a posteriori numerical certificate}

A dense numerical eigensolve costs $O(n^3)$ arithmetic and $O(n^2)$ storage. After preparation, a dense implementation constructs the final direction in $O(nk)$ arithmetic; its coefficients must be communicated or programmed. The reference implementation also evaluates its compressed Rayleigh term densely. In the ideal paired basis that term can be precomputed from the unselected component of $b$, reducing decoding work. None of this removes the physical work of performing the linear functional.

For \emph{any} prepared orthonormal first basis and any unit final direction $w$ in its complement, the exact conditional-first-stage risk of the correct conditional predictor is
\[
 2\left(\norm{QAQ}_F^2-(w^TAw)^2\right)
       +\norm t^2-(t^Tw)^2. \tag{18}
\]
The predictor must use $(t^Tw)z$, not $\norm t\,z$ when $w$ is only approximately aligned. Equation (18) exposes both residual quadratic energy and lost linear alignment. If $w$ is obtained from an approximate vector $\widetilde t\in\operatorname{range}(Q)$, the latter loss is at most $\norm{t-\widetilde t}^2$. If the implemented decoder instead uses $\norm{\widetilde t}\,z$, its exact risk is the quadratic term in (18) plus $\norm{t-\widetilde t}^2$. In the noisy $QAQ=0$ construction, the corresponding excess risk is $\alpha\norm{t-\widetilde t}^2$. These absolute-error certificates remain useful near $t=0$. A skipped last read pays $\norm t^2$ in the exact setting.

Approximate diagonalization need not assume an eigengap for an approximate-risk statement. Rational plane rotations can remain exactly orthogonal while driving the off-diagonal Frobenius residual below a chosen tolerance. The read Astra N1164 source explicitly constructs such gap-free rational Jacobi rotations for known matrices \cite{astra1164}; its unknown-source acquisition remains separate. A real-symmetric scaling/shift gives the corresponding diagonalization operation here. Treat the diagonalized nearby matrix as $\widehat A$, and pay its off-diagonal error through (16)--(17). Exact inertia classification, in contrast, requires certified signs.

Rational approximations to the pairing angles can leave residual paired Rayleigh values $d_i$. If the rotations remain within certified orthogonal eigenpair blocks, their extra compressed energy is exactly $2\sum_i d_i^2$. With a general approximate eigenbasis, off-block terms must also be counted; (18) provides the general certificate. An implemented orthogonality or model residual must be bounded, not assumed away.

For a known correlated Gaussian $X=\mu+LZ$, apply the construction to
$A_Z=L^TAL$ and $b_Z=L^T(2A\mu+b)$. Acquiring/factoring the covariance and implementing the transformed measurements are additional costs. Small covariance eigenvalues can amplify physical measurement norms and noise. An unknown or merely approximately Gaussian law is not covered automatically.

\subsection{What the end-to-end fixture actually paid}

The reference implementation receives only scalar access to a rotated, dense, 32-dimensional quadratic with a linear term. It pays 561 model evaluations and 17,952 submitted model coordinates, with a supplied absolute source-error bound $10^{-6}$. The acquired matrix has inertia $(16,16,0)$. It then pays the eigensolve and uses 16 first-batch readings plus one final reading per deployment.

The actual acquired-model squared $L^2$ error was $1.420\times10^{-9}$; the conservative bound (16) was $1.2449\times10^{-8}$. Across 8,000 deployments with exact latent reads, observed MSE was $1.407\times10^{-9}$. Across a separate 8,000 deployments with noise variance 0.01 on every read, observed MSE was $0.020326$, with Monte Carlo standard error $0.000333$; the acquired-model risk formula was $0.020465$. Each run paid 136,000 scalar reads and 4,352,000 dense inner-product coordinate contributions.

These are synthetic consistency executions. The source error bound and quadratic promise were generated by the fixture; they were not inferred from a real device. The floating eigensolve has recorded numerical residuals, rather than a formal interval-arithmetic certificate.

If a full evaluator remains live, an ordinary alternative is to acquire the full input and call it once, avoiding offline model reconstruction. If coefficients or a calibrated model are already available, every comparator receives them. Learning costs may dominate a small number of deployments. Reuse can amortize preparation, but no end-to-end runtime superiority follows from these examples.

\section{Coordinate-only companion: condition away interactions}

The spectral result does not retroactively grant linear measurements to the previous coordinate model. A separate construction works at the original access interface.

Let independent coordinates have arbitrary known laws and a square-integrable target represented by node and pairwise factors on a graph $G$. For any vertex cover $H$, first read $H$. Conditional on those values, the remaining target is a sum of independent one-coordinate functions. Read the $k$ remaining coordinates with largest conditional variance contributions.

\begin{theorem}[Vertex-cover two-round comparison]
This policy has risk no greater than every original hard-$k$ adaptive coordinate policy, using at most $\min(n,|H|+k)$ reads. Only second moments are required.
\end{theorem}
\begin{proof}
Given $H$, no interaction remains between unobserved coordinates. Their individual variance contributions do not change when other remaining coordinates are observed, so the largest $k$ are optimal even against adaptive continuations. Reveal $H$ for free and emulate any original policy, supplying already known values internally. Its remaining reads number at most $k$ and its information includes the original transcript. The explicit second batch dominates this emulator at each first-stage history.
\end{proof}

For an explicit sparse graph, a maximal matching produces a cover within twice the minimum size in $O(n+|E|)$ work. Standard bipartite methods can find a minimum cover when applicable. Conditional coefficient updates and variance selection are ordinary sparse scans; prior moments and tables remain acquired inputs.

On the preserved rare-gate family with $m$ gates and $2m-1$ leaves per gate, the gates form a size-$m$ cover. Reading them and then up to $2m$ active leaves uses at most $3m$ reads, dominates every hard-$2m$ adaptive policy, and has risk tending to $e^{-1}$. Every fixed $3m$-read design has risk tending to one. At $m=20$, an exact ordinary allocation calculation gives risk $0.351713$ for two rounds and $0.925641$ for the optimal fixed 60-read design. This addresses final accuracy at the same coordinate interface.

The cheap rule is a hard-budget result. Under an expected-budget comparator, favorable first-stage outcomes may justify buying many more leaves. A global threshold on conditional variance can solve that allocation, but integrating over first-stage outcomes may be expensive. The companion proof preserves both the failure of the naive expected-budget transfer and its separately priced extension.

\section{Grouped model acquisition and its useful failed fixture}

A separate companion reduces the previous coordinate-by-coordinate influence acquisition under an exact promised Boolean quadratic oracle. Four value contrasts give nonnegative unbiased estimators for a group's total influence mass and for omitted-component risk. Their bounded relative variance supports median-of-means tests. A best-first splitting tree either acquires an approximately ranked set or certifies that the remaining maximum influence is small relative to residual risk.

Its proved scalar-call bound is polynomial in $k,1/\varepsilon,\log n,\log(1/\alpha)$, rather than linear in $n$ merely to scan every coordinate at every context. The displayed construction uses $O((k/\varepsilon^2)\log n)$ group/coverage tests, each with a paid confidence sample count, plus retained-coefficient recovery. A residual mean estimate replaces exact full-intercept recovery when a relative-risk allowance is available. This is a separate Boolean access model, not an unpriced learner of arbitrary spectral directions.

The executed small fixture deliberately exposes the current constants: it paid 1,140,791 scalar calls, while full ordinary interpolation required 137. This is a theoretical acquisition construction with a failed efficiency fixture. Dense input traffic and source execution multiply the scalar count. The exact quadratic promise remains necessary; a noisy influence floor does not establish the end-to-end multiplicative guarantee for arbitrarily small residual risk. The complete proof, code, strong junta-learning comparisons and Astra N1071 grouping ancestry are in the companion files.

\section{Evidence, primary comparisons and remaining limits}

The central independent audit completed 384 tensor Gauss--Hermite risk checks, 192 compression checks, 576 arbitrary-projection lower-bound checks and 480 pointwise reconstructions. The conditional squared residual is degree four, so the three-point Gaussian rule is exact in exact arithmetic; these executions are floating point. The maximum relative risk discrepancy was $3.51\times10^{-15}$.

The noisy audit completed 27 conditional-risk quadratures, 360 arbitrary/repeated fixed-design checks, 90 optimum-attainment checks and 21 rational tangent checks. A separate 40,000-input Monte Carlo run per method at $m=128,s=0.01$ gave MSEs $0.019764$, $0.507788$ and $0.019643$ for two rounds, optimal fixed 129 reads and fixed 256 reads, respectively. The corresponding standard errors were $0.000141$, $0.003648$ and $0.000139$. These checks support formulas and implementation; the proofs establish the general statements.

The coordinate companion used exact rational Bellman and allocation calculations. Reviews also retained the model-stencil blindness example, the same-budget threshold improvement, the high-noise repeated-read counterexample, and the distinction between hard and expected exact-query complexity.

\begin{longtable}{>{\raggedright\arraybackslash}p{.25\textwidth}>{\raggedright\arraybackslash}p{.66\textwidth}}
\toprule Primary machinery & Relationship to this construction\\\midrule\endhead
Classical Gaussian information complexity \cite{ww,ritter,jackowski} & Wasilkowski--Wo\'zniakowski's primary abstract states a no-adaptation result for linear problems under orthogonally invariant measures. Jackowski's 1991 average-case multilinear paper is a closer antecedent: its publisher record was retrieved, but its exact cross-argument adaptive-information contract was not. Historical priority of the combined theorem remains unresolved.\\[4pt]
Adaptive sensing limits \cite{acd} & Sparse-vector recovery under noisy norm-constrained observations is an important established comparison. Estimating one nonlinear functional has a different objective from recovering the vector. Our fixed-design lower bound is proved for the actual functional and interface.\\[4pt]
Quadratic goal-oriented design \cite{goal} & The Gaussian quadratic variance and data-averaged fixed-design objective are established in Theorems 3.1--3.2. Their preselected sensor placements, computational approximations and matrix-free model access are substantive alternatives; the present joint adaptive comparison is a different information question.\\[4pt]
Ordinary spectral approximation \cite{fanpall} & Leading eigenvectors supply the optimal homogeneous fixed exact experiment (9). The signed compression minimum follows from Fan--Pall's classical theorem; maximal isotropic subspaces are also classical. Any conventional method may adopt the paired compiler.\\[4pt]
Few-round stochastic optimization \cite{ghuge,gns} & Existing few-round and adaptivity-gap results use stated stochastic cover/probing objectives. They motivate the comparison, but do not substitute for the conditional-risk proofs here.\\[4pt]
Existing learning and grouping \cite{ei,astra1071} & Low-degree junta learning and positive group queries are material acquisition ancestors. Our Boolean companion claims a specific risk-competitive construction, not a general learning-speed breakthrough.\\\bottomrule
\end{longtable}

The important completed distinction is between \emph{feedback depth} and \emph{number of observations}. For the declared quadratic linear-measurement problem, unrestricted sequential planning can be replaced by one feedback round with one additional exact read. For a concrete noisy functional, the retained feedback produces a large difference in final error against the optimal fixed experiment. At coordinate-only access, a separately acquired interaction cover supplies another two-round reduction.

The remaining scientific burdens are substantial: establish an important application with the required linear measurement interface; acquire its model and noise law economically; pay direction control and precision; and compare against live full-evaluator, task-specific and ordinary spectral alternatives in total cost. General noisy adaptive optimality, arbitrary correlated/non-Gaussian laws, unknown high-order targets and historical priority are not established by this report.

\clearpage
\begin{thebibliography}{20}
\bibitem{ww} G. W. Wasilkowski and H. Wo\'zniakowski (1984), \emph{Can adaption help on the average?}, Numerische Mathematik 44, 169--190. Primary author-institution abstract: \url{https://scholars.uky.edu/en/publications/can-adaption-help-on-the-average/}. DOI: \url{https://doi.org/10.1007/BF01410103}.
\bibitem{ritter} K. Ritter, \emph{Nonlinear methods for linear problems}, chapter in Average-Case Analysis of Numerical Problems. Public publisher references read; full body unavailable. \url{https://link.springer.com/chapter/10.1007/BFb0103941}. Corresponding nonlinear-problem chapter: \url{https://link.springer.com/chapter/10.1007/BFb0103942}.
\bibitem{jackowski} T. Jackowski (1991), \emph{Complexity of multilinear problems in the average case setting}, Journal of Complexity. Publisher metadata verified; full information contract unavailable through attempted public routes. \url{https://doi.org/10.1016/0885-064X(91)90047-2}.
\bibitem{fanpall} K. Fan and G. Pall (1957), \emph{Imbedding conditions for Hermitian and normal matrices}, Canadian Journal of Mathematics 9, 298--304. Theorem 1 and the real-symmetric corollary at the end of Section 2 read in the primary publisher PDF. \url{https://doi.org/10.4153/CJM-1957-036-1}.
\bibitem{goal} J. N. Neuberger, A. Alexanderian and B. van Bloemen Waanders (2024 preprint), \emph{Goal oriented optimal design of infinite-dimensional Bayesian inverse problems using quadratic approximations}. Primary HTML read, especially Theorems 3.1--3.2 and Sections 1--4. \url{https://arxiv.org/html/2411.07532v1}.
\bibitem{acd} E. Arias-Castro, E. J. Cand\`es and M. A. Davenport, \emph{On the Fundamental Limits of Adaptive Sensing}. \url{https://arxiv.org/abs/1111.4646}.
\bibitem{ghuge} R. Ghuge, A. Gupta and V. Nagarajan, \emph{The Power of Adaptivity for Stochastic Submodular Cover}, ICML 2021. Primary abstract read; no unretrieved constants imported. \url{https://proceedings.mlr.press/v139/ghuge21a.html}.
\bibitem{gns} A. Gupta, V. Nagarajan and S. Singla, \emph{Adaptivity Gaps for Stochastic Probing: Submodular and XOS Functions}. \url{https://arxiv.org/abs/1608.00673}.
\bibitem{ei} A. Eskenazis and P. Ivanisvili, \emph{Learning low-degree functions from a logarithmic number of random queries}. \url{https://arxiv.org/html/2109.10162v3}.
\bibitem{astra1164} Astra Research Knowledge, N1164 and its source \emph{Finite positive-cost rounding for one-sided effects}, Section 2. The complete SHA-256 and pinned source path are in the evidence ledger. Known-matrix rational Jacobi construction; acquisition and physical execution are not supplied by that lemma.
\bibitem{astra1071} Astra Research Knowledge, N1071: sparse causal graphs and factorial quotients. Positive group-query and zero-versus-gap premises retained in the companion comparison.
\end{thebibliography}
The repository navigation, packets and governing source ranges were read at pinned revision
\begin{center}\small\texttt{2fd26b68f3e2dff7520167aa05f23ee1054337e3}.\end{center}
Start route: \href{https://github.com/harjassand/astra-research-knowledge/blob/2fd26b68f3e2dff7520167aa05f23ee1054337e3/00_START_HERE.txt}{\texttt{astra-research-knowledge / 00\_START\_HERE.txt} at the pinned revision}.
Capital integration remains separate; no repository changes were pushed.
\end{document}
